\section{First interactions and indexed singular values}\label{tra:first-interactions}

Coordinate deletion tracks a representation statistic through a split. The next argument instead stops paths at their first interactions. We combine its sparse norm estimate with a projection overlap bound to retain the singular rank inside each irreducible block.

This method controls every singular index within an irreducible representation. Its sparse argument orders the first interactions of labels and deletes them in reverse order. Its dense argument charges coloring constraints to conditional total correlation. The resulting rank estimate closes under a balanced split of the bit directions.

Write $T_n$ for the sweep operator, $V_\lambda$ for an irreducible
representation of $S_n$, and $D_\lambda=\dim V_\lambda$. Each trace is
unnormalized. The index in the next theorem belongs to a single
irreducible block; its contribution to a regular trace is repeated
$D_\lambda$ times. We first control blocks close to the trivial
representation directly from their paths. A compatibility estimate
then handles the other blocks and permits induction on the bit length.

\begin{theorem}\label{tra:first-interactions:singular}

\label{tra:first-interactions:indexed-conclusion}
There is an absolute finite $p_*>0$ such that, for every dyadic
$n$, every partition $\lambda\vdash n$, and every $1\le i\le D_\lambda$,
\begin{equation}\label{tra:first-interactions:eq1}
 s_i(T_n(\lambda))\le(D_\lambda i)^{-1/(2p_*)}.
\end{equation}
The singular values are listed in decreasing order within one copy of
$V_\lambda$.

\end{theorem}

\subsection{A sparse estimate by deleting first interactions}

A representation whose first row has length $n-k$ first occurs in
the action on $k$ ordered distinct cards. On its new part, any term
that ignores a card vanishes. This cancellation forces every tracked
card to interact with another one. Our immediate task is to show
that such complete coverage has small operator norm when
$k\le n^{.99}$. The following occupation estimate will control the
exceptional starting configurations.

\begin{lemma}[Occupation bounds for partially coupled paths]
\label{grid:occupation}
Start a set of $r$ cards at distinct sites of the binary cube and move
them by one complete coordinate sweep. Independently, move $s$ further
walkers, each by its own fair bit in every coordinate; these walkers
may have coincident endpoints. For disjoint sets of sites
$A_1,\ldots,A_L$, let $N_i$ count all $r+s$ endpoints in $A_i$, with
repetitions counted for the independent walkers. For nonnegative
integers $h_i$,
\[
 \mathbb E\prod_i\binom{N_i}{h_i}
 \le\prod_i\frac{((r+s)|A_i|/n)^{h_i}}{h_i!}.
\]
In particular the same right side bounds
$\Pr\{N_i\ge h_i\text{ for every }i\}$. The assertion also holds
with the coordinate order reversed.
\end{lemma}
\begin{proof}
Let $\eta(x)$ be the occupation indicator for the $r$ jointly moved
cards, and $p_x=\mathbb E\eta(x)$. At the deterministic initial set,
$\mathbb E\prod_{x\in A}\eta(x)\le\prod_{x\in A}p_x$ for every set
$A$. A fair switch on $u,v$ preserves this family of inequalities.
If $A$ contains neither endpoint there is no change. If it contains
one endpoint, averaging the two old inequalities replaces its marginal
by $(p_u+p_v)/2$. If it contains both, the occupation product is
unchanged and $p_up_v\le((p_u+p_v)/2)^2$. Disjoint switches can be
applied successively, proving preservation through the sweep. At its
end every $p_x=r/n$, because the endpoint of each individual card is
uniform. Summing over distinct tested sites in the disjoint $A_i$
therefore bounds the joint binomial moments for these cards by
$\prod_i(r|A_i|/n)^{h_i}/h_i!$.

The independent walkers have multinomial factorial moments bounded
by the same expression with $s$ in place of $r$: assign the tested
hits to distinct walkers, and use their uniform independent endpoints.
Expand each binomial coefficient of the sum of the two counts using
Vandermonde's identity. Independence of the two populations and the
binomial theorem replace $r$ and $s$ by $r+s$. This proves the stated
bound. No part of this proof depends on which coordinate is updated
first.
\end{proof}

\begin{lemma}\label{tra:first-interactions:sparse}

\label{tra:first-interactions:sparse-statement}
For some absolute \(c_s>0\) and all sufficiently large \(n\), if \(1\le k=n-\lambda_1\le n^{.99}\), then
\begin{equation}
                        \|T_n(\lambda)\|_{\rm op}\le n^{-c_s k}.\label{tra:first-interactions:eq2}
\end{equation}

\end{lemma}\begin{proof}\label{tra:first-interactions:sparse-proof}

\label{tra:first-interactions:endpoint-kernel-and-centered-cover}
To see this, take the sweep kernel on \(k\)-tuples, with row variable \(x=(x_1,\ldots,x_k)\) the input tuple and column variable \(y\) the output tuple. Actions on tuples, or on functions of them, estimate the same singular values on \(\lambda\). For two tuple labels \(u,v\), set
\[
 h=h(u,v)=\max\{i:(x_u)_i\ne(x_v)_i\},\qquad
 l=l(u,v)=\min\{i:(y_u)_i\ne(y_v)_i\}
\]
for distinct tuple entries as here. Define
\[
                   w(h,l)=\begin{cases}1& l<h,\\
                                        2& l=h,\\
                                        0& l>h.
                           \end{cases}
\]
Then the tuple kernel is
\begin{equation}
                         n^{-k}\prod_{u<v} w(h(u,v),l(u,v)).\label{tra:first-interactions:eq3}
\end{equation}
Indeed given endpoints, paths replace one bit after another in \(x\) by the corresponding bit of \(y\). They must remain nonoverlapping as paths through positions. For a pair, if the output bits below \(h\) (indices \(<h\)) agree, the paths would use the same switch at \(h\), where output bits must be opposite. When compatible they use either no common switch or that one, by the bit replacement order. Thus pair conditions suffice, and each shared switch gains a factor 2 compared to probabilities for independent paths.

For \(J\subset [k]\), put
\[
           Q_J(x,y)=n^{-k}\prod_{\{u,v\}\subset J} w(h(u,v),l(u,v)).
\]
This is the restriction (to our distinct tuples) of the kernel using the joint shuffle only on labels in \(J\), and placing the other output labels independently uniformly; likewise its transpose is restricted from using the reverse shuffle jointly only on \(J\). In particular each of its row and column sums is at most 1. In bounding the block on \(\lambda\) we can replace \eqref{tra:first-interactions:eq3} by
\[
                         Z=\sum_{J\subset[k]}(-1)^{k-|J|}Q_J .
\]
Indeed proper-subset terms can be dropped on compression to the \(\lambda\)-isotypic space. Their images or the images of their transposes are in subspaces of functions depending only on a shorter tuple, which have no \(\lambda\) part.

Make the interaction graph on \([k]\) for \((x,y)\) by using an edge for \(u,v\) if \((y_u)_i=(y_v)_i\) for all \(i<h(u,v)\); the edge time is \(h(u,v)\). The sum for \(Z\) vanishes if there is an isolated label, by cancellation. Thus it suffices for norm bounds to use nonnegative kernels \(Q_J\), summed but restricted by the coverage condition (all labels have an interaction).

\label{tra:first-interactions:bad-cell-factorial-moments}
Put \(\epsilon=.01\), so that \(k\le n^{1-\epsilon}\), and take \(\alpha=\epsilon/2\). We use the dyadic cells from the input hierarchy: a cell of size \(2^h\) specifies the bits of index \(>h\). For an input \(x\), call a label bad if it is in some such cell of size \(t\) where the label count \(r\) from \(x\) satisfies
\[
                            r\ge 2,\qquad r> t n^{-\alpha}.
\]
Otherwise call the label good. Split rows into bad rows (at least \(k/2\) bad labels) and the other rows (called good rows).

Here are two probability bounds for this split.

\begin{itemize}
\item Given \(y\), the column sum of \(Q_J\) over bad rows is \(n^{-\Omega(k)}\), meaning bounded by \(n^{-c k}\) for an absolute \(c>0\) for all sufficiently large \(n\), uniformly. In the transpose experiment the output considered here is \(x\). We will use factorial moment bounds on counts in cells. Lemma~\ref{grid:occupation}, with $r=|J|$ and $s=k-|J|$,
gives, for disjoint cells of sizes $t_i$ and required counts $r_i$,
\begin{equation}\label{tra:first-interactions:eq4}
 \Pr\{N_i\ge r_i\text{ for all }i\}
 \le\prod_i\frac{(kt_i/n)^{r_i}}{r_i!}.
\end{equation}
This bound is used before restriction to distinct output tuples;
that restriction can only decrease a column sum of the nonnegative
kernel.

  A bad row gives a family of disjoint cells (take maximal ones from the bad-label definition), where
  \[
                  r_i\ge 2,\qquad t_i/r_i\le n^\alpha,\qquad
                           k/2\le R_0=\sum r_i\le k
  \]
  with \(r_i\) their counts. We spell out a union bound. For choices of ordered sizes and counts, multiplying \eqref{tra:first-interactions:eq4} by \(\prod_i n/t_i\) allows the choices of locations. Using \(r!\ge(r/e)^r\), the result is at most
  \[
        \prod_i \left[e^{r_i}\frac{k}{r_i}
                   \left(\frac{k t_i}{n r_i}\right)^{r_i-1}\right]
               \le e^k k^L n^{-(\epsilon-\alpha)(R_0-L)}.
  \]
  In summing over ordered lists we divide by \(L!\), since we seek unordered families of distinct cells. The factor \(k^L/L!\) is at most \(e^k\). The remaining number of size and count lists over lengths and totals is bounded by \((C(d+1))^{Ck}\) using compositions, for some absolute \(C\). Since \(R_0-L\ge R_0/2\ge k/4\), this proves the column bound.
\end{itemize}

\label{tra:first-interactions:first-meeting-deletion}
\begin{itemize}
\item Given \(x\) in a good row, the row sum of \(Q_J\) subject to coverage is \(n^{-\Omega(k)}\), uniformly. Simulate \(Q_J\) before restriction by using independent personal fair bits for each label at each time. Labels outside \(J\) just use these bits; labels in \(J\) also do, except that when two in \(J\) use the same switch we use one of their bits (chosen, say, by label index) and its opposite for the two respective new position bits. This generates the kernel since placements of the \(J\) labels need only this joint use of bits, independently at each switch. Only \(J\)-labels count for this bit-sharing rule.

  We track the first interaction of a label in the interaction graph (now defined by the output bits whether or not all outputs end up distinct). Until its first interaction its output bits come from its own personal bits. Indeed an actual \(J\)-pair at a switch at time \(h\) must have highest differing input bit \(h\) and matching lower output bits, as the \(J\)-positions have stayed distinct.

  If all good labels have interactions, we can select a subset \(D\) of them of size at least half the number of good labels, and for each \(u\in D\) specify a partner \(v(u)\) witnessing its first interaction, so that if the partner also belongs to \(D\), its first time is strictly earlier. To justify the selection, look at labels having their first interaction at a given switch, defined by its time \(h\), the fixed higher input bits, and the common lower output bits. On either of the two sides of this switch there is at most one label having its first interaction there. Two labels on that same side, arriving with the same lower output bits, would have an interaction at their own highest differing input bit, which would be earlier. Hence at most two have their first interaction there. In case two good labels first interact there, we can omit one to avoid ties; one good label alone with first interaction there needs no omission (partners not required to be selected).

  For fixed specifications with times \(h_u=h(u,v(u))\), the probability of these witnesses with the stated first-time property is bounded by
  \begin{equation}
                                   \prod_{u\in D} 2^{-(h_u-1)}.\label{tra:first-interactions:eq5}
  \end{equation}
  To see this without independence assumptions about the partners, take labels in \(D\) of largest specified time \(h\), and compare with the simulation deleting those labels, using the same personal bits of the rest. On the witness event the rest follow the deletion simulation through time \(h-1\), since no deleted label has interacted before then. Earlier witnessing events persist in that simulation. Each deleted label must match its specified partner on the first \(h-1\) output bits, and on the event uses personal bits up to then. Its partner belongs to the rest by the selection condition. With the rest's personal bits given, these necessary matches have probability \(2^{-(h-1)}\) per deleted label, independently for the purpose of this bound. This gives \eqref{tra:first-interactions:eq5} by successive deletion; equivalently we sum specifications only where within \(D\) partners are of strictly earlier specified time.

  For one good label \(u\), the sum over possible partners of \(2^{-(h(u,v)-1)}\) is at most \(2d n^{-\alpha}\). For a contributing level \(h\), the cell of size \(2^h\) containing \(u\) contains a partner and so contains at most \(2^h n^{-\alpha}\) labels since \(u\) is good. Because in a good row there must be a witness specification with \(|D|\ge k/4\) if coverage holds, summing \eqref{tra:first-interactions:eq5} over the subsets and partners gives the row bound for large \(n\).

\end{itemize}
For \(Z\) we can sum these bounds with a \(2^k\) factor. On the good-row portion, one has small row sums in absolute value by the second bound and at most \(2^k\) for absolute column sums. On the bad-row portion, the first bound gives small column sums, and row sums in absolute value are at most \(2^k\). The row-column sum bound on operator norm now gives \(\|Z\|_{\rm op}\le n^{-\Omega(k)}\). This proves \eqref{tra:first-interactions:eq2}.

\end{proof}

\subsection{Conditional total correlation in a rectangle}

\label{tra:first-interactions:balanced-grid-setup}
The sparse argument is complete: ignoring any one label kills the
new representation level, and the remaining path systems are rare on
good rows or on bad columns. We now turn to the large-dimensional
blocks. A balanced split reduces them to products of smaller sweep
operators. The missing input is a bound on the overlap of their
Fourier subspaces, which we obtain by estimating the probability of
row--column compatibility. In it we regard positions as an \(a\times b\) rectangle, \(a b=n\), with
\[
                       a,b\in[\sqrt n/\sqrt 2,\sqrt 2\sqrt n].
\]
Use subgroups \(H,J_0\) of the symmetric group \(G\) on the rectangle: \(H\) permutes positions separately within each column (so \(H=(S_a)^b\) in symmetric group notation), and \(J_0\) does so in each row. Below, logs for entropy and estimates are natural. We continue to allow all bound constants to be absolute and to consider sufficiently large sizes.

We will need the following probability estimate.

\begin{proposition}\label{tra:first-interactions:grid}

\label{tra:first-interactions:grid-statement}
Draw elements \(h\in H\) (column permutations \(h_v,\ 1\le v\le b\)) and \(j\in J_0\) (row permutations \(j_r,\ 1\le r\le a\)) with all permutations independent, from possibly nonuniform laws. Suppose on each factor the pointwise density versus uniform is bounded, with the product of those upper bounds \(e^z\). Then
\begin{equation}
              \Pr\{ h j\in J_0 H\}
                  \le \exp\{-n+C n^{.98} + C z/\log n\}.\label{tra:first-interactions:eq6}
\end{equation}
Here \(h j\) means apply \(j\) and then \(h\). The coefficient of $-n$ is essential: the later regular-trace
normalization contributes $+n$ to leading order.

\end{proposition}\begin{proof}\label{tra:first-interactions:grid-proof}

\label{tra:first-interactions:total-correlation-exposure}
For given \(j\), make a bipartite multigraph between left vertices \(u\) and right vertices \(v\), both sets indexing columns. Each initial site \((r,u)\) gives an edge, of row \(r\), from \(u\) to \(v=j_r(u)\). Each vertex has degree \(a\). Color this edge by \(h_v(r)\), a color from \([a]\), so colors are all used once at each right vertex. If \(hj\in J_0 H\), every left vertex also sees every color once, since cards of each initial column must go one to each final row. Call this a valid coloring. (This condition also suffices, though necessity is enough.)

Condition the experiment on the event in \eqref{tra:first-interactions:eq6}, assuming positive probability. Its log inverse probability is the relative entropy of this new law with respect to the original experiment. Write it as
\[
                   L=L_j+L_h+L_{\rm corr},
\]
where:
\begin{itemize}
\item \(L_j\) is the relative entropy paid for the new marginal of \(j\);
\item \(L_h\) is the sum, averaged over this marginal of \(j\), of relative entropies paid for the individual column permutation marginals conditional on \(j\), versus their original independent laws;
\item \(L_{\rm corr}\) is the conditional total correlation among the column permutations given \(j\), i.e., the relative entropy against the product of their conditional marginals (averaged over \(j\)).

\end{itemize}
All are nonnegative. Write \(z_j,z_h\) for the parts of \(z\) from the row, column factors, respectively. The relative entropy of the new marginal of \(j\) versus independent uniform row permutations is at most \(L_j+z_j\). The sum of relative entropies of the column permutation marginals given \(j\), each versus uniform, is in expectation at most \(L_h+z_h\). These comparisons follow from the density upper bounds.

Estimate \(L_{\rm corr}\) by exposing right vertices one by one. In this argument we are in the conditioned experiment, in particular the coloring is valid. Fix \(j\), and take an order from independent uniform priorities in \([0,1]\) for the right vertices, independent of everything else. By the chain rule, the conditional total correlation is the sum of the mutual informations between each \(h_v\) and the preceding \(h\)-variables (for fixed order, given \(j\)). Within \(h_v\) expose entries in row order \(1,\ldots,a\) and use the chain rule again. Consider the entry corresponding to a particular edge, from left vertex \(u\) to \(v\), and condition also on earlier entries at \(v\), but not yet other columns' information. Denote by \(p_c\), \(c\in[a]\), the conditional law of this entry. We use mutual information of the color with the preceding vertices, for the conditional law (and average over these conditionings). For a fixed set of preceding vertices, if the color is \(c\), none of those vertices can have color \(c\) on an edge from \(u\). Thus, writing mutual information as the average relative entropy of the preceding vertices' variables given the color versus their marginal here, a lower bound for this term is
\begin{equation}
             \sum_c p_c \left(-\log\Pr\{\text{color }c
                           \text{ not used from }u\text{ at preceding vertices}\}\right).\label{tra:first-interactions:eq7}
\end{equation}
This uses the marginal not given the color inside the \(-\log\), with the conditioning on \(j\) and earlier entries at \(v\) still in force; the bound follows from relative entropy for a law supported on an event.

If \(v\)'s priority is \(t\), average the preceding set over the other priorities. For color \(c\), write
\[
  \delta_c=\Pr\{\text{the edge from }u\text{ colored }c\text{ goes to }v\}
\]
in the same conditional law not given the entry. Validity implies that the survival probability inside \eqref{tra:first-interactions:eq7}, averaged over priorities at the other vertices, is \(1-t+t\delta_c\). Jensen's inequality and integrating in \(t\) now lower bound the expected term by
\begin{equation}
                           \sum_c p_c \int_0^1 -\log(1-t+t\delta_c)\,dt.\label{tra:first-interactions:eq8}
\end{equation}
For \(\delta_c=0\) the integral is 1. For all \(0\le s\le1\), the integral with \(\delta_c=s\) is nonnegative and at least \(1-C\sqrt s\): the loss from 1 in the interior is \(s\log(1/s)/(1-s)\).

\label{tra:first-interactions:parallelism-and-heavy-atom-losses}
Call an edge heavy in parallelism if its endpoints have edge multiplicity greater than \(a^{1/2}\). We may discard \eqref{tra:first-interactions:eq8} for these edges. On another edge \(\sum_c\delta_c\le a^{1/2}\). For the colors with \(p_c\le a^{-.7}\), we have
\[
             \sum_{c:p_c\le a^{-.7}} p_c\sqrt{\delta_c}
                  \ \le\ \Big(\sum_{c:p_c\le a^{-.7}} p_c\delta_c\Big)^{1/2}
                  \ \le\ a^{-.1}.
\]
For the other colors we allow a loss up to their total \(p_c\)-mass instead in \eqref{tra:first-interactions:eq8}. Consequently, summed over edges and averaged, \eqref{tra:first-interactions:eq8} bounds \(L_{\rm corr}\) below by \(n\), less the expected number of edges heavy in parallelism, the sum of the exceptional masses with \(p_c>a^{-.7}\) (averaged), and \(C n a^{-.1}\).

We check the sizes of these losses using the marginal divergences above. We use the elementary estimate
\begin{equation}
                D(P\Vert Q)\ \ge\ s P(F)-\log(1+Q(F)(e^s-1))\label{tra:first-interactions:eq9}
\end{equation}
for \(s>0\) and an event \(F\), by the entropy inequality (or just projecting to \(F\)). For \(h_v\) given \(j\), in the chain of exposing entries use as reference a uniform permutation. Except at most \(\lceil a^{.9}\rceil\) late entries per column, it has at least \(a^{.9}\) remaining possibilities. The uniform conditional chance to hit the exceptional set for the current \(p\) is then at most \(a^{-.2}\) (there are at most \(a^{.7}\) such colors). Use \eqref{tra:first-interactions:eq9} with \(s=.09\log a\) at these exposures, with the set fixed conditional on earlier entries and \(j\). The relative entropies sum by the chain rule. This bounds the total expected exceptional mass over the edges, including the late-entry loss, by
\begin{equation}
                     C\big((L_h+z_h)/\log a+n a^{-.1}\big).\label{tra:first-interactions:eq10}
\end{equation}

For parallelism, consider the new marginal law of \(j\) against the reference of independent uniform permutations. Expose \(j\) by rows (all of one row permutation, then the next) and within each by successive entries. For an entry corresponding to left vertex \(u\), count a hit if the target vertex has occurred as target from \(u\) already at least \(a^{.4}\) times before the entry. There are at most \(a^{.6}\) possible hit targets given the earlier entries. At least half the edges heavy in parallelism give hits (in the sense that their endpoint pairs with multiplicity above \(a^{1/2}\) give hits on at least half those occurrences, for large \(a\)). Again discount at most \(\lceil b^{.9}\rceil\) late exposures per permutation where the number of possibilities under uniform may be small. Otherwise the uniform conditional chance of a hit is at most \(C b^{-.3}\), since \(a,b\) are comparable. Applying \eqref{tra:first-interactions:eq9} with \(s=.09\log b\) as before, the expected number of edges heavy in parallelism is at most
\begin{equation}
                      C\big((L_j+z_j)/\log b+n b^{-.1}\big).\label{tra:first-interactions:eq11}
\end{equation}
This bounds the loss even if \(j\) acquires strong correlations upon conditioning.

The lower bound from \eqref{tra:first-interactions:eq8}, with \eqref{tra:first-interactions:eq10} and \eqref{tra:first-interactions:eq11}, gives
\[
                     L_{\rm corr}\ge n-C n^{.98}
                                   - C\,(L_h+L_j+z)/\log n
\]
where we use comparability and have loosened the numerical power. In \(L\), the losses involving \(L_h,L_j\) are paid for by those terms themselves. Thus for sufficiently large \(n\) this proves \eqref{tra:first-interactions:eq6}.

\end{proof}

\subsection{Crossing projections and dyadic singular bands}

\label{tra:first-interactions:subgroup-product-setup}
Here is how we use the rectangle estimate. By splitting the bit list into segments of lengths \(\lfloor d/2\rfloor,\lceil d/2\rceil\), the sweep product takes the form
\[
                                A B ,
\]
where \(A,B\) are group-algebra operators on \(H,J_0\), respectively. The roles of rows and columns may be named to have this form. Each operator on its own subgroup is a tensor product (law of independent sweeps) on the separate columns or rows, with the corresponding smaller sizes \(a,b\). All estimates below are unchanged if using adjoint operators consistently.

\begin{lemma}\label{tra:first-interactions:crossing}

\label{tra:first-interactions:projection-fourth-moment}
For orthogonal projections \(P,Q\) in the subgroup algebras of \(H,J_0\), suppose each is a tensor product over that subgroup's factors. Let their regular ranks, meaning their ranks in the regular representation of the respective subgroup, be \(r_H,r_J\), assumed nonzero. Put \(z=\log(r_H r_J)\). We claim that on any irreducible \(\lambda\) of \(G\),
\begin{equation}
          D_\lambda \| (P Q)(\lambda) \|_4^4
                 \ \le\ \exp\big(C n^{.98}+(1+C/\log n)z\big),\label{tra:first-interactions:eq12}
\end{equation}
where \(\|\cdot\|_4\) is Schatten 4-norm.

\end{lemma}\begin{proof}

\label{tra:first-interactions:regular-coefficient-proof}
It suffices to bound the full regular trace of $PQPQ$, since its
irreducible contributions are nonnegative fourth powers with regular
multiplicity. Lemma~\ref{lem:coefficient-compatibility}, applied to
the local projections, gives
\begin{equation}
 \Tr_{\rm reg}(PQPQ)
 \le\frac{|G|}{|H||J_0|}\,r_Hr_J\,
       \Pr\{hj\in J_0H\}.
 \label{tra:first-interactions:eq13}
\end{equation}
Here the line permutations are independent under their normalized
squared-coefficient laws, each with density bounded by its local
regular rank. These ranks multiply to $r_Hr_J=e^z$, so
\eqref{tra:first-interactions:eq6} applies with exactly the chosen $z$.
Finally
\[
                    \frac{|G|}{|H|\,|J_0|}
                    =\frac{n!}{(a!)^b (b!)^a}
                    \ \le\ \exp(n+C\sqrt n\log n)
\]
by Stirling's formula. Combining in \eqref{tra:first-interactions:eq13} proves \eqref{tra:first-interactions:eq12}, allowing adjustment of the constant in \(C n^{.98}\). The main exponential here is why we used the coloring estimate in addition to ranks.

\end{proof}

\begin{proof}[Proof of Theorem~\ref{tra:first-interactions:singular}]

\label{tra:first-interactions:alt-and-dyadic-band-decay}
We perform the singular-value induction. The preceding fourth-moment estimate has an error independent of
the eventual Schatten exponent. We retain that property while summing
the spectral bands; it lets the induction begin from any sufficiently
large finite exponent without changing the cutoff in $n$. Suppose for the two smaller sizes the bound in \eqref{tra:first-interactions:eq1} holds with \(p_0\) (in place of \(p_*\)), where we can take the larger of the two parameters, always at least 2. We will show how along induction it suffices to increase this parameter by factors at most \(1+O(1/\log n)\) for all sufficiently large sizes.

Take
\[
                    p=p_0(1+C_0/\log n)
\]
with \(C_0\) large enough, absolute. For each \(\lambda\), we will bound a Schatten trace of \(A B\) on the irreducible. Write \(X=A^* A\), \(Y=B B^*\). Lemma~\ref{lem:positive-power-comparison}, applied on the
irreducible block with $A_2=A$, $A_1=B$, $P=2p$ and $q=4$, gives
\begin{equation}
                  \operatorname{Tr}|A B|^{2p}
                      \ \le\ \| X^{p/4} Y^{p/4}\|_4^4 .\label{tra:first-interactions:eq14}
\end{equation}

Decompose in the subgroup algebras to bound the right side. For \(X\), on each of the \(b\) column factors take the irreducible matrix-block decomposition and, within each irreducible, the eigenbasis of the corresponding factor of \(X\), in order of decreasing eigenvalue. Group indices within this matrix-block into dyadic intervals
\([R,2R)\) truncated at its dimension (here \(R=1,2,4,\ldots\)). Each such interval gives a projection in that factor group algebra of regular rank
\[
                              m \cdot (\text{length of interval}),
\]
where \(m\) is the dimension of this smaller irreducible. Taking tensor products gives projections \(P\) as in \eqref{tra:first-interactions:eq12} for \(H\) that sum to identity. By the smaller-size hypothesis, for each such choice
\begin{equation}
                          \|P X^{p/4}\|_{\rm op}
                                   \le r_H^{-p/(4p_0)},\label{tra:first-interactions:eq15}
\end{equation}
where the norm bound holds also after lifting to \(G\). Indeed squared singular values from an interval on a factor are at most \((m R)^{-1/p_0}\); the interval length is at most \(R\). The analogous decomposition of \(Y\) over row factors gives \(Q\) with \(\|Q Y^{p/4}\|_{\rm op}\le r_J^{-p/(4p_0)}\).

\label{tra:first-interactions:band-count-and-intermediate-trace}
The total number of pairs \(P,Q\) in these decompositions is at most
\begin{equation}
                                      \exp(C n^{.8}).\label{tra:first-interactions:eq16}
\end{equation}
In fact per factor of size \(a\) or \(b\), the number of partitions for its irreducibles is at most \(\exp(C\sqrt a)\) or \(\exp(C\sqrt b)\), and the number of dyadic intervals in each is at most polynomial in the factor size since dimensions are at most the factorial. The partition bound here is the standard elementary one; for example it follows by using \(s=1/\sqrt a\) and bounding the partition generating function at \(e^{-s}\) by \(\exp(C/s)\), using its product formula (and likewise with \(b\)).

Use the triangle inequality for Schatten 4-norm to sum terms on the right of \eqref{tra:first-interactions:eq14}, inserting these projections. For each \(P,Q\), by \eqref{tra:first-interactions:eq12}, \eqref{tra:first-interactions:eq15} and the ideal property of Schatten norm, we have, on \(\lambda\),
\[
 \begin{aligned}
 D_\lambda \| X^{p/4} P Q Y^{p/4} \|_4^4
   &\le \exp\big(- (p/p_0)z + C n^{.98}+(1+C/\log n)z\big)\\
   &\le \exp(C n^{.98}),
 \end{aligned}
\]
with our choice of \(C_0\). Now summing by \eqref{tra:first-interactions:eq16} (triangle inequality before taking fourth powers) still gives by \eqref{tra:first-interactions:eq14}
\begin{equation}
                         D_\lambda \operatorname{Tr}|T_n(\lambda)|^{2p}
                                  \le \exp(C_1 n^{.98}),\label{tra:first-interactions:eq17}
\end{equation}
where \(C_1\) is absolute. The placement of the Schatten powers in this proof means we do not have a loss from summing norm bounds raised later to a possibly very large \(p\).

\label{tra:first-interactions:dimension-and-all-index-closure}
To use \eqref{tra:first-interactions:eq17} and finish the induction, we record simple dimension bounds. For shapes of height exceeding \(n/2\) the singular bound was automatic. For the rest with \(k=n-\lambda_1\ge n^{.99}\), we have
\begin{equation}
                              \log D_\lambda\ge c n^{.99}\label{tra:first-interactions:eq18}
\end{equation}
for large \(n\), with an absolute \(c>0\). Here and later we can see this type of lower bound as follows. Transpose a partition if needed and let \(u\) be its first row length after taking the larger of width and height. If \(u<n/8\), the hook lengths are at most \(2u\), so the dimension is at least \((n/(2eu))^n\) by the hook formula. For \(u\ge n/8\), put \(t=n-u\). The hook formula gives at least
\begin{equation}
                 \binom{n}{t}\exp\left(-C\frac{t}{u}\log(u+1)\right),\label{tra:first-interactions:eq19}
\end{equation}
since we may drop the dimension (a factor at least 1) of the partition below that row, and the correction to row factorial uses log factors
\(\log(1+c_j/(u-j+1))\), where \(c_j\) counts cells below in column \(j\). Indeed $jc_j\le t$, and
\[
 \sum_{j=1}^{u}\log\left(1+\frac{c_j}{u-j+1}\right)
 \le t\sum_{j=1}^{u}\frac{1}{j(u-j+1)}
 =\frac{2t}{u+1}\sum_{j=1}^{u}\frac1j.
\]
This proves the displayed correction, including the range in which
the tail is not small. For \(n/8\le u\le n/2\), \eqref{tra:first-interactions:eq19} gives log dimension of order at least \(c' n\). For \(u>n/2\) under our height restriction before transposition, \(u=\lambda_1\), so \(t=k\); \eqref{tra:first-interactions:eq18} follows using \eqref{tra:first-interactions:eq19} and \(\binom{n}{t}\ge (n/t)^t\).

Hence in this dense case \eqref{tra:first-interactions:eq17} yields
\[
 \begin{aligned}
                s_i(T_n(\lambda))
                    &\le \left(\frac{\exp(C_1 n^{.98})}{D_\lambda i}\right)^{1/(2p)}
                     \le (D_\lambda i)^{-1/(2p')},
 \end{aligned}
\]
where we can take \(p'=p(1+C_2/\log n)\) for a suitable absolute \(C_2\) and sufficiently large \(n\), by \eqref{tra:first-interactions:eq18}.

For \(1\le k\le n^{.99}\), \eqref{tra:first-interactions:eq2} suffices, with a fixed absolute induction parameter, since \(D_\lambda i\le n^{2k}\) using occurrence on tuples and \(i\le D_\lambda\). And for \(k=0\) the induction bound just says \(1\le 1\). This completes the bounds used in the induction step.

\label{tra:first-interactions:bounded-induction-parameter}
To formalize uniformity, use a fixed threshold for sufficiently large sizes, and below that threshold take a common finite induction parameter making the singular inequality true. This is possible by the strict finite-size gap in Lemma~\ref{lem:finitegap}. Enlarge this starting parameter if necessary to be at least 2 and to handle the fixed parameter from \eqref{tra:first-interactions:eq2}. For \(n=2^d\) above threshold, by the step just proved the desired parameter need only multiply the larger child parameter (children of bit lengths \(\lfloor d/2\rfloor,\lceil d/2\rceil\)) by at most \(1+C/d\), changing absolute constants. Their product is uniformly bounded by Lemma~\ref{lem:dyadic-exponents} with $\gamma=1$. Thus an absolute \(p_*\) works for all sizes. This proves \eqref{tra:first-interactions:eq1}.

\end{proof}

\subsection{Mixing after a bounded number of sweeps}

\label{tra:first-interactions:fourier-summation-and-time-normalization}
The estimate \eqref{tra:first-interactions:eq1} is an indexed
bound on each irreducible block, with exponent $a=1/(2p_*)$ in
the notation of Section~\ref{sec:spectral-conversion}. Its final
conversion gives
\[
 4\|\mu_q-U\|_{\TV}^2
 \le\sum_{\lambda\ne(n),(1^n)}D_\lambda^{\,2-q/p_*}
 \longrightarrow0
\]
for any fixed integer $q$ with $q/p_*\ge3$; sign and all taller
annihilated shapes contribute zero. This uses matrix powers through
Schatten or operator norm bounds and does not require normality of a
sweep. Hence an absolute number of sweeps gives vanishing error and
$t_{\mathrm{mix}}(d)=O(d)$. Lemma~\ref{lem:support} supplies the
physical-time lower bound $2d-O(1)$.
